\section{The canonical distance and its fibres}\label{app:collapsed}

This appendix begins an alternative proof of the nef-and-big application,
Corollary~\ref{thm:nef-big}. It uses the canonical metric and the geometry of
its collapsed fibres, rather than the holomorphic box-extension theorem
used in the proof of Theorem~\ref{thm:main}. The three subsequent
appendices construct crossings around each whole fibre, prove that these
crossings do not increase the canonical distance in either variable,
and use that estimate to fill finite grids along arbitrary paths.

Let \(X\) satisfy the hypotheses of Corollary~\ref{thm:nef-big}.
Proposition~\ref{canon:contraction} gives the crepant canonical morphism
\(f:X\to Z\), and Lemma~\ref{canon:metric-input} gives its limiting
K\"ahler--Einstein metric \(g\) on
\(X^0=f^{-1}(Z_{\reg})\). By Proposition~\ref{canon:parallel}, the
supplied summands are parallel on \(X^0\) and integrable on all of
\(X\). We first make them orthogonal without changing the completion,
then establish the two topological properties of the fibres that the
crossing argument needs.

\begin{proposition}\label{canon:pseudodistance}
There is a real-analytic Kähler metric \(h\) on \(X^0\), uniformly
equivalent to \(g\), for which the given splitting is orthogonal and
locally a Riemannian product.  Its intrinsic distance completion is
canonically homeomorphic to \(Z\).  Denoting the resulting distance
on \(Z\) by \(\rho\), the function
\begin{equation}\label{canon:d}
 d(x,y)=\rho(f(x),f(y))
\end{equation}
is a continuous pseudodistance on \(X\), and \(d(x,y)=0\) exactly
when \(f(x)=f(y)\).
\end{proposition}

\begin{proof}
Let \(Q_j:T_X\to E_j\) denote projection along the other supplied
summand, and define
\[
 h(v,w)=g(Q_1v,Q_1w)+g(Q_2v,Q_2w).
\]
Proposition~\ref{canon:parallel} implies
\(\nabla^gQ_j=0\), hence \(\nabla^gh=0\).
Thus \(\nabla^g\) is also the Levi-Civita connection of \(h\).
It preserves the complex structure, so \(h\) is Kähler.

The positive \(g\)-self-adjoint endomorphism representing \(h\)
is parallel.  Its eigenvalues are constant on connected \(X^0\).
Consequently there are global positive constants \(a,b\) such that
\[
 ag\leq h\leq bg.
\]
In local product coordinates, first- and second-factor real coordinate
fields \(V,W\) commute.  The equality
\(\nabla_VW=\nabla_WV\), together with preservation of the
complementary summands, makes both sides zero.  Since \(\nabla h=0\),
the first block of \(h\) depends only on first coordinates and the
second block only on second coordinates.

A Kähler--Einstein metric is real analytic in holomorphic coordinates
and their conjugates \cite[Theorem~6.1(c)]{DK81}.  The \(Q_j\) are
holomorphic, so the displayed formula for \(h\) proves its real
analyticity.

Uniform equivalence of the metrics makes the identity on \(X^0\)
bi-Lipschitz for their intrinsic distances.  It therefore extends
uniquely to a bi-Lipschitz homeomorphism of their completions.
Lemma~\ref{canon:metric-input} identifies that completion with \(Z\),
extending the given identification on \(X^0\).  The assertions about
\(d\) follow from continuity of \(f\) and from \(\rho\) being a
distance.
\end{proof}

\begin{lemma}\label{canon:fibers}
For every \(z\in Z\), the reduced fiber
\(P=f^{-1}(z)_{\mathrm{red}}\) is a compact connected projective
algebraic set.  Any two points of \(P\) are joined inside \(P\) by a
chain of images of rational curves, and \(\pi_1(P)=1\) in the ordinary
complex topology.
\end{lemma}

\begin{proof}
Projectivity and connectedness follow from the construction of \(f\).
The rational-chain assertion is the fiber theorem for resolutions of
klt varieties \cite[Corollary~1.6, arXiv v2]{HM07}.

For simple connectedness we use the resolution-fiber consequence of
Takayama's local argument: the fiber of a proper birational morphism
from a smooth complex variety to a klt variety is simply connected in
the ordinary topology
\cite[Theorem~1.1 and the first paragraph on p.~829]{Takayama03}.
This consequence is also used explicitly in
\cite[proof of Theorem~3.10, arXiv v4]{NS24}.
Here \(X\) is smooth and \(Z\) is klt by
Proposition~\ref{canon:contraction}, so it applies to \(f:X\to Z\)
and gives \(\pi_1(f^{-1}(z))=1\).
Replacing the fiber by its reduction does not change its underlying
topological space, hence \(\pi_1(P)=1\).
\end{proof}

The two conclusions of Lemma~\ref{canon:fibers} have distinct roles.
Rational chains will supply continuation routes, whereas ordinary simple
connectedness will remove their monodromy. The latter conclusion is not
inferred from rational chain connectedness of a possibly singular fibre.
